\section{Pairwise Preference-Residual Diagnostics on Ant}
\label{app:ant_diagnostics}

Complementing the candidate-search ablations in Table~\ref{tab:ant_search}, we
tested whether explicit nonlinear preference-interaction capacity improves
interior execution once actor and critic updates are aligned to the queried
preference $u$. Across five independent seeds at 75k interactions
(seeds 1301--1305), we crossed a base own-preference actor against a
pairwise-residual own-preference actor under both progress-routed and
quota-regulated portfolio allocation, alongside native CAPQL. The residual
actor augments the pre-tanh action mean with zero-initialized bilinear
preference heads:
\begin{equation}
 \mu(s, w, z) = \mu_{\mathrm{base}}(h)
       + \sum_{1 \leq i < j \leq m} 2 w_i w_j (A_{ij} h + b_{ij}),
 \label{eq:pairwise_residual}
\end{equation}
where $h$ denotes the shared trunk activation. Because $w_i w_j = 0$ at all
one-hot simplex vertices, the residual branch activates exclusively on mixed
preferences while sharing trunk representations. Outcomes are evaluated over
the 14 mode-0 weights satisfying $0 < w_{\mathrm{energy}} < 1$
(Table~\ref{tab:ant_pairwise}).

% Generated by make_evidence_tables.py from completed readouts.
\begin{table}[t]
\centering\small
\caption{Ant structure-by-selector pilot on five independent seeds at 75k. Base and residual BAMOR arms both use own-preference actor and critic updates. $\Delta U_0$ is mixed-energy utility in mode 0; HV averages all modes.}
\label{tab:ant_pairwise}
\begin{tabular}{lrrrr}
\toprule
Contrast & $\Delta U_0$ & W/L/T & Final HV & W/L/T \\
\midrule
Residual effect, adaptive & +39.6 & 2/3/0 & +1.62\% & 3/2/0 \\
Residual effect, blind & -412.5 & 1/4/0 & +0.27\% & 3/2/0 \\
Selector effect, residual & +177.9 & 3/2/0 & +0.70\% & 4/1/0 \\
Selector effect, base & -274.2 & 0/5/0 & -0.64\% & 2/3/0 \\
Residual adaptive vs CAPQL & -247.9 & 2/3/0 & -2.06\% & 1/4/0 \\
\bottomrule
\end{tabular}
\end{table}


Under progress-routed allocation, adding pairwise residual heads yields a modest
mean gain of $+39.6$ mixed-energy utility (2/5 wins). Under score-neutral
allocation the same change gives $-412.5$. The resulting difference in
differences is positive on all five seeds (mean $+452.1$). With the base
actor, progress-based routing lost to the score-neutral control on all five
seeds ($-274.2$ mixed-energy utility). The progress-routed residual actor
still trailed native CAPQL by $247.9$ mixed-energy utility and $2.06\%$ Final HV
($25$ runs totaling $1{,}875{,}000$ training and $4{,}700{,}168$ evaluation
interactions). To preserve a single task-agnostic architecture across all
benchmarks, the bilinear head was kept strictly diagnostic and not incorporated
into the main HalfCheetah and Hopper evaluations.
